\documentclass[hyphens,aspectratio=169,dvipsnames]{beamer}
\usepackage{graphicx}
\usepackage{xcolor}
\usepackage{amssymb}
\usepackage{amsmath}
\usepackage{minted}
\usetheme{Berlin}
\usecolortheme[RGB={0,95,47}]{structure}
\beamertemplatenavigationsymbolsempty

\makeatletter
\AtBeginEnvironment{minted}{\dontdofcolorbox}
\def\dontdofcolorbox{\renewcommand\fcolorbox[4][]{##4}}
\makeatother

\begin{document}

\begin{frame}[fragile]{\texttt{curl}}
    \begin{itemize}
        \pause \item \texttt{curl} is a command-line tool for speaking various network protocols.
        \pause \item We have it installed on plink.
        \pause \item Example!
    \end{itemize}
\end{frame}

\begin{frame}[fragile]{HTTP}
    \begin{itemize}
        \pause \item \texttt{curl} speaks many protocols, but nearest to my heart is HTTP :)
        \pause \item You can use \texttt{curl} to make HTTP requests by giving it a URL on \texttt{argv}.
        \pause \item \texttt{curl http://example.com}
    \end{itemize}
\end{frame}

\begin{frame}[fragile]{What is a URL?}
    \begin{itemize}
        \pause \item What does a URL do?
        \pause \item What are the pieces of a URL? (Whiteboard)
    \end{itemize}
\end{frame}

\end{document}
